% !TEX program = xelatex
\documentclass[a4paper,12pt]{article}

\usepackage{fontspec}
\usepackage{polyglossia}
\setdefaultlanguage{russian}
\setmainfont{Times New Roman}

\usepackage{geometry}
\geometry{top=1.8cm,bottom=1.8cm,left=2.3cm,right=2.3cm}

\usepackage{enumitem}
\usepackage{titlesec}

% Компактные заголовки
\titlespacing*{\section}{0pt}{0.8em}{0.4em}

% Компактные списки
\setlist[itemize]{leftmargin=*,itemsep=0.6em,topsep=0.4em,parsep=0pt}

\title{\textbf{Список тем курсовых работ и проектов}}
\author{}
\date{}

\begin{document}
\maketitle

\section*{Темы проектов}

\begin{itemize}

\item \textbf{Гардероб (Крутова).}
Приложение для офлайн-хранения гардероба пользователя и подбора сочетаемых элементов одежды.

\item \textbf{Отметка посещаемости (Крутова).}
Наряду с мобильным приложением требуется бэкенд (например, FastAPI). Приложение должно обеспечивать удобную отметку посещаемости, аналитику для преподавателя, передачу данных через API. Необходимо продумать загрузку групп и списков, вопросы безопасности и соблюдения законодательства.

\item \textbf{Зоократия (или иные настольные игры, ранее не оцифрованные).}
Использование возможно только в учебных целях из-за лицензионных ограничений (при отсутствии договорённости с правообладателями). Требуется реализация настольной игры в виде мобильного приложения с возможным взаимодействием через бэкенд или по Bluetooth.

\item \textbf{Поиск пути на велосипеде с учётом ПДД + навигация.}
Существующие приложения некорректно учитывают требования ПДД РФ. Предполагается офлайн-работа с картами (например, OpenStreetMap). Ожидается прототип.

\item \textbf{Поиск пути на велосипеде с учётом ПДД: исследование алгоритмов.}
Исследование алгоритмов поиска пути с учётом ограниченных ресурсов смартфона. Возможна командная работа: один участник реализует приложение, другой — различные алгоритмы и их тестирование.

\item \textbf{Поиск пути на велосипеде с указанием мест ДТП.}
Использование существующих решений с дополнительным получением данных о ДТП через публичные API и их визуализацией. Регион — Москва.
Возможна командная работа с авторами предыдущих тем.

\item \textbf{Велосипед и общественный транспорт: поиск маршрута.}
Офлайн-поиск маршрутов с велосипедными и транспортными сегментами. Учитываются варианты использования собственного и прокатного велосипеда, а также ограничения перевозки. Регион — Москва.
Возможна командная работа с авторами предыдущих тем.

\item \textbf{Средства индивидуальной мобильности и общественный транспорт.}
Аналог предыдущей темы с использованием СИМ вместо велосипеда.

\item \textbf{Маршруты наземного общественного транспорта с учётом типа маршрута.}
Офлайн-поиск маршрутов с отображением категории: магистральный, социальный, районный. Тип маршрута учитывается при минимизации времени ожидания. Регион — Москва.
Возможна командная работа с авторами предыдущих тем.

\item \textbf{Поиск маршрута по критерию минимальной неопределённости.}
Оптимизация по дисперсии времени в пути, а не по среднему значению. Предпочтение отдаётся более предсказуемым маршрутам.
Возможна командная работа с авторами предыдущих тем.

\item \textbf{Маршрут пешком и на общественном транспорте для маломобильного пользователя.}
Учитываются ограничения инфраструктуры: станции метро и МЦД, лестницы, уклоны, доступность остановок.
Возможна командная работа с авторами предыдущих тем.

\item \textbf{Отслеживание фитнес-активности (сенсоры и видео).}
Автоматическое определение количества и качества упражнений (например: отжимания, приседания, подтягивания) с использованием минимум двух механизмов анализа — сенсорного и видео. Формируются аналитические отчёты.

\item \textbf{Автозапись перемещения на общественном транспорте (ожидание и стоимость).}
Фоновое определение поездок на автобусе, электробусе, трамвае, метро и пригородных поездах с фиксацией первой и последней остановки, времени ожидания и расчётом затрат при различных способах оплаты в Москве.

\item \textbf{Автозапись перемещения (анализ поведенческих паттернов).}
Анализ типичных сценариев использования транспорта по сезонам и дням недели. Время ожидания не является ключевым параметром.

\item \textbf{Автозапись перемещения (анализ скорости и проблемных участков).}
Определение маршрутной скорости на автобусах, электробусах и трамваях с визуализацией замедлений и аналитикой.

\item \textbf{Исследование энергопотребления фонового трекинга.}
Сравнение различных способов фонового определения местоположения и оптимизация по критерию энергопотребления.

\item \textbf{Определение способа передвижения.}
Классификация перемещения: пешком, велосипед, СИМ, автобус, электробус, трамвай, метро. Анализ точности и энергоэффективности методов.

\item \textbf{Маршрут для прогулки с заданными критериями.}
Построение офлайн-маршрута по параметрам: зелёные зоны, отсутствие лестниц, интересные здания, туристические точки.

\item \textbf{Анализ инфраструктуры с позиции велосипедиста.}
Оценка выбранной области с точки зрения доступности и качества велосипедной инфраструктуры.

\item \textbf{Офлайн-поиск маршрутов с пересадками (межгород).}
Построение маршрутов на поездах, автобусах и электричках с учётом пересадок (возможно использование Yandex API).

\end{itemize}

\end{document}
